\section{Related Work}

\subsection{Autoformalization and Repair}

LLMs have been explored to reduce the effort of translating informal natural language requirements into formal specifications~\cite{zhang2025position}.
One direction is property synthesis, where natural language is mapped to LTL~\cite{cosler2023nl2spec, fuggitti2023nl2ltl}, though this covers only a fragment of specification languages and cannot capture full system models.
Other work addresses program specification generation~\cite{wen2024enchanting, wu2023lemur} or proof synthesis and repair~\cite{first2023baldur, lu2024proof, carrott2024coqpyt}, but these remain task-specific rather than providing an integrated modeling and verification workflow.

More recent studies target system-level models, such as Alloy~\cite{hong2025effectiveness} and Event-B~\cite{capozucca2025ai}, but these approaches are limited in scope and rely on direct text-to-code mappings.
The closest work is PAT-Agent~\cite{zuo2025pat}, which synthesizes PAT models and validates them via model checking.
While effective for complex systems and requirements, its verification is confined to full-state-space model checking.
By contrast, Event-B Agent combines bounded model checking with theorem proving, enabling counterexample-guided repair and proof-based correctness.

% Beyond LLM-based work, the B and Event-B community has studied complementary directions such as construction and refinement methods~\cite{alkhammash2015building,dupont2021event,mashkoor2017refinement}, automated model repair~\cite{cai2019automatic,cai2022fast,kobayashi2024repairing}, and code generation frameworks for producing executable implementations~\cite{mery2011automatic,furst2014code,rivera2017code}. In comparison, prior B or Event-B efforts have focused on human-guided construction, specialized repair techniques, or code generation, whereas Event-B Agent brings these strands together under an automated and unified framework.
Beyond LLMs, the B and Event-B community has advanced construction and refinement methods~\cite{alkhammash2015building,dupont2021event,mashkoor2017refinement}, model repair~\cite{cai2019automatic,cai2022fast,kobayashi2024repairing}, and code generation~\cite{mery2011automatic,furst2014code,rivera2017code}.
These efforts remain human-guided, repair-specific, or downstream, whereas Event-B Agent unifies synthesis, repair, and verification in an automated framework.

% PAT-Agent: \cite{zuo2025pat}
% Alloy Synthesis: \cite{hong2025effectiveness}
% Event-B Synthesis: \cite{capozucca2025ai}
% LTL Synthesis: \cite{cosler2023nl2spec, fuggitti2023nl2ltl}
% Program Specification Synthesis: \cite{wen2024enchanting, wu2023lemur}
% Proof Synthesis: \cite{first2023baldur, lu2024proof, carrott2024coqpyt}

\subsection{Neurosymbolic Methods}
Neurosymbolic methods \cite{chaudhuri2021neurosymbolic} combine the expressiveness of neural learning with the verifiability of symbolic reasoning.
Recent directions include minimal neural activation patterns for verifiable robustness~\cite{geng2024learning}, automata embeddings for reinforcement learning with semantic guarantees~\cite{yalcinkaya2025provably}, and specification learning from combined membership and preference queries~\cite{shah2023learning}.
Other efforts integrate logic into differentiable frameworks, such as tensorized LTL\textsubscript{f} constraints verified in Isabelle/HOL and utilized for neural training~\cite{chevallier2025formally}, or world models that follow principles for physical interpretability, such as structuring latent spaces by physical intent and partitioning outputs for verifiability~\cite{peper2025four}.

Event-B Agent shares this neurosymbolic philosophy but targets formal verification, using LLMs for semantic tasks while ensuring correctness through model checking and theorem proving in a unified synthesis-and-repair loop.

\subsection{Proof Tactic Recommendation}
% Recommending tactics in interactive theorem proving has long been studied to lower users' burden. Early studies focused on learning from human proofs and replaying strategies in a given prover. 
% Examples include PaMpeR~\cite{nagashima2018pamper}, which provides proof method recommendations for Isabelle/HOL, TacticToe~\cite{gauthier2021tactictoe}, which learns and searches tactic sequences in HOL4 using Monte Carlo Tree Search, and Tactician~\cite{blaauwbroek2020tactician}, which records tactic applications in Coq and reuses them to suggest or search for proofs. 
% PSL~\cite{nagashima2017proof} further introduced a high-level proof-strategy language that expands into concrete tactic sequences during execution.

Tactic recommendation in interactive theorem proving has long aimed to reduce user effort. Early approaches learned from existing proofs and replayed strategies, including PaMpeR for Isabelle/HOL~\cite{nagashima2018pamper}, TacticToe for HOL4~\cite{gauthier2021tactictoe}, Tactician for Coq~\cite{blaauwbroek2020tactician}, and PSL, a strategy language that expands into concrete tactic sequences~\cite{nagashima2017proof}.

% With the rise of transformers and large language models, recent research has shifted toward retrieval- and LLM-based approaches. 
% MagnusHammer~\cite{mikula2023magnushammer} improves premise selection using transformer-based contrastive learning, while LeanDojo~\cite{yang2023leandojo} builds a retrieval-augmented prover that combines LLMs with large proof libraries. 
% Other directions explore proof-step prediction for Lean using fine-tuned LLMs~\cite{gloeckle2023temperature}, recursive proof generation strategies~\cite{wang2024proving}, and lifelong or curriculum-based learning agents such as LeanAgent~\cite{kumarappan2024leanagent}. 
% Lean-STaR~\cite{lin2024lean} further integrates chain-of-thought style reasoning with tactic prediction, interleaving planning and proving to improve success rates.
Recent work leverages transformers and LLMs. MagnusHammer~\cite{mikula2023magnushammer} enhances premise selection with contrastive learning, while LeanDojo~\cite{yang2023leandojo} combines LLMs with large proof libraries in a retrieval-augmented prover. Other directions explore proof-step prediction~\cite{gloeckle2023temperature}, recursive proof generation~\cite{wang2024proving}, lifelong agents such as LeanAgent~\cite{kumarappan2024leanagent}, and approaches such as Lean-STaR~\cite{lin2024lean}, which interleave chain-of-thought reasoning with tactic prediction to improve proof success.

These efforts center on tactic prediction and retrieval within provers, whereas Event-B Agent embeds tactic guidance in a broader neurosymbolic framework, where model construction, verification, and repair progress together rather than as isolated tasks.

% \subsection{Event-B Synthesis and Repair}

% \textcolor{red}{TODO: add this only if there's enough space.\\
% TODO: neurosymbolic methods?}